\documentclass[11pt]{article}

\usepackage[a4paper,margin=1in]{geometry}
\usepackage[T1]{fontenc}
\usepackage[utf8]{inputenc}
\usepackage{lmodern}
\usepackage{graphicx}
\usepackage{amsmath}
\usepackage{enumitem}
\usepackage{hyperref}
\usepackage{xcolor}
\usepackage{float}
\usepackage{booktabs}
\usepackage{array}

\hypersetup{
    colorlinks=true,
    linkcolor=blue!50!black,
    urlcolor=blue!50!black
}

\setlist[itemize]{leftmargin=1.5em, itemsep=0.3em, topsep=0.3em}

\title{Time-Varying Break-Even Growth Rates for the Swiss Labour Market:\\Line Manager Brief}
\author{}
\date{\today}

\begin{document}
\maketitle

\section{Executive Summary}
\begin{itemize}
    \item Benchmark result: in the direct \(Y \rightarrow U\) specification, the implied break-even GDP growth rate is about \(1.9\%\) annualized on average and about \(2.1\%\) annualized in the latest quarter, 2025 Q4.
    \item Time variation is present but modest under the preferred calibration with \(\sigma_{\eta}=0.05\,\sigma_{OLS}\): the direct \(Y \rightarrow U\) break-even series has a standard deviation of about \(0.2\) percentage points annualized.
    \item Additional model results:
    the \(Y \rightarrow E\) specification implies a negative break-even GDP growth rate on average,
    while the \(E \rightarrow U\) specification implies break-even employment growth of about \(0.9\%\) annualized on average and \(1.2\%\) annualized in 2025 Q4.
    \item Internal consistency check: chaining \(Y \rightarrow E\) and \(E \rightarrow U\) gives an implied \(Y \rightarrow U\) break-even rate close to the direct estimate in levels, but more volatile over time.
    \item Practical read: the direct \(Y \rightarrow U\) specification should remain the main benchmark; the employment-based specifications are useful companion views and interpretation aids.
\end{itemize}

\begin{table}[H]
\centering
\small
\begin{tabular}{>{\raggedright\arraybackslash}p{3.2cm} >{\raggedleft\arraybackslash}p{2.6cm} >{\raggedleft\arraybackslash}p{2.9cm} >{\raggedleft\arraybackslash}p{2.8cm}}
\toprule
Model & Break-even threshold & Scenario 1: \(Y=0\%\) / \(E=0\%\) & Scenario 2: \(Y=2\%\) / \(E=1\%\) \\
\midrule
\(Y \rightarrow U\) & \(2.1\) & \(0.1\) & \(0.0\) \\
\(Y \rightarrow E\) & \(-1.2\) & \(0.4\) & \(1.1\) \\
\(E \rightarrow U\) & \(1.2\) & \(0.1\) & \(0.0\) \\
\bottomrule
\end{tabular}
\caption{Latest model readings for 2025 Q4. Growth rates are annualized.}
\end{table}

\section{Question and Motivation}
\begin{itemize}
    \item Objective: estimate the pace of Swiss real GDP growth needed to stabilize labour-market conditions, and allow that threshold to vary over time.
    \item Policy relevance: this is a compact way to summarize how labour-market resilience may be changing even if output growth, unemployment, and employment all remain noisy in the short run.
    \item Strategy: start from transparent reduced-form relationships, then let only the intercept vary over time in a local-level state-space model.
\end{itemize}

\begin{figure}[H]
    \centering
    \includegraphics[width=0.98\textwidth]{../outputs/memo_motivation_scatters.png}
    \caption{Motivation from the raw reduced-form relationships among GDP growth, employment growth, and unemployment changes.}
\end{figure}

\section{Time-Invariant Evidence}
\begin{itemize}
    \item Baseline static reduced-form equation:
    \[
    U_t = \alpha + \beta_0 Y_t + \beta_1 Y_{t-1} + \beta_2 Y_{t-2} + \varepsilon_t.
    \]
    \item Estimated benchmark equation:
    \[
    U_t = 0.123 - 0.108\,Y_t - 0.092\,Y_{t-1} - 0.070\,Y_{t-2} + \varepsilon_t,
    \]
    with \(R^2 \approx 0.48\).
    \item Interpreting ``break-even'' as zero unemployment change and assuming constant GDP growth across current and lagged terms implies:
    \[
    Y_t^\ast = -\frac{\alpha}{\beta_0+\beta_1+\beta_2}.
    \]
    \item Static benchmark:
    \(Y_t^\ast \approx 0.455\%\) q/q, or about \(1.83\%\) annualized.
    \item ARDL robustness check:
    \[
    U_t = \alpha + \phi U_{t-1} + \beta_0 Y_t + \beta_1 Y_{t-1} + \beta_2 Y_{t-2} + \varepsilon_t.
    \]
    \item Estimated ARDL equation:
    \[
    U_t = 0.054 + 0.714\,U_{t-1} - 0.105\,Y_t - 0.015\,Y_{t-1} - 0.007\,Y_{t-2} + \varepsilon_t,
    \]
    with \(R^2 \approx 0.80\) and much better residual behavior.
    \item The implied ARDL break-even GDP growth is close to the simpler benchmark:
    about \(0.431\%\) q/q, or \(1.74\%\) annualized.
\end{itemize}

\section{Main Time-Varying Results}
\begin{itemize}
    \item We use the same local-level state-space routine across the three reduced-form relationships, but only illustrate the mechanics in detail for the benchmark \(Y \rightarrow U\) specification.
    \item Preferred benchmark specification:
    \[
    U_t = c_t + \beta_0 Y_t + \beta_1 Y_{t-1} + \beta_2 Y_{t-2} + \gamma_2 D_{2020Q2,t} + \gamma_3 D_{2020Q3,t} + \varepsilon_t,
    \]
    with
    \[
    c_t = c_{t-1} + \eta_t.
    \]
    \item Estimation is done in \texttt{statsmodels} with a local-level model; the state innovation standard deviation is calibrated to \(\sigma_{\eta}=0.05\,\sigma_{OLS}\).
    \item The implied break-even GDP growth rate in the benchmark model is:
    \[
    Y_t^\ast = -\frac{c_t}{\beta_0+\beta_1+\beta_2}.
    \]
    \item Main findings:
    average break-even GDP growth is \(1.86\%\) annualized,
    the latest reading in 2025 Q4 is \(2.09\%\) annualized,
    and the annualized standard deviation is only about \(0.20\) percentage points.
    \item Scenario interpretation:
    under GDP stagnation, the model implies unemployment rises by about \(0.12\) percentage points on average;
    under GDP growth of \(2\%\) annualized, the implied unemployment change is essentially zero on average and \(0.006\) percentage points in 2025 Q4.
    \item Subsection 1, \(Y \rightarrow E\):
    \[
    E_t = c_t + \beta_0 Y_t + \beta_1 Y_{t-1} + \beta_2 Y_{t-2} + \gamma_2 D_{2020Q2,t} + \gamma_3 D_{2020Q3,t} + \varepsilon_t.
    \]
    This maps GDP growth into employment growth rather than unemployment changes.
    It implies an average break-even GDP growth rate of \(-0.76\%\) annualized and a latest reading of \(-1.18\%\) annualized in 2025 Q4; under GDP growth of \(2\%\) annualized, implied employment growth is about \(0.98\%\) annualized on average and \(1.14\%\) annualized in 2025 Q4.
    \item Subsection 2, \(E \rightarrow U\):
    \[
    U_t = c_t + \theta_0 E_t + \theta_1 E_{t-1} + \theta_2 E_{t-2} + \gamma_2 D_{2020Q2,t} + \gamma_3 D_{2020Q3,t} + \varepsilon_t.
    \]
    This yields a break-even employment-growth concept,
    \[
    E_t^\ast = -\frac{c_t}{\theta_0+\theta_1+\theta_2}.
    \]
    Average break-even employment growth is \(0.91\%\) annualized and the latest reading is \(1.16\%\) annualized in 2025 Q4; with employment growth of \(1\%\) annualized, implied unemployment change is close to zero on average and \(0.017\) percentage points in 2025 Q4.
    \item Bottom line across the three specifications:
    the direct \(Y \rightarrow U\) model remains the most interpretable headline measure,
    while \(Y \rightarrow E\) and \(E \rightarrow U\) add complementary information on employment dynamics and labour-market transmission.
\end{itemize}

\begin{figure}[H]
    \centering
    \includegraphics[width=0.83\textwidth]{../outputs/memo_y_to_u.png}
    \caption{Preferred benchmark: time-varying \(Y \rightarrow U\) specification. Top: implied break-even GDP growth. Bottom: implied unemployment change under two GDP scenarios.}
\end{figure}

\begin{figure}[H]
    \centering
    \includegraphics[width=0.83\textwidth]{../outputs/memo_y_to_e.png}
    \caption{Time-varying \(Y \rightarrow E\) specification. Top: implied break-even GDP growth for employment. Bottom: implied employment growth under two GDP scenarios.}
\end{figure}

\begin{figure}[H]
    \centering
    \includegraphics[width=0.83\textwidth]{../outputs/memo_u_on_e.png}
    \caption{Time-varying \(E \rightarrow U\) specification. Top: implied break-even employment growth. Bottom: implied unemployment change under two employment-growth scenarios.}
\end{figure}

\section{Direct Versus Chained \(Y \rightarrow U\)}
\begin{itemize}
    \item A useful internal check is to combine the two employment-based steps:
    first \(Y \rightarrow E\), then \(E \rightarrow U\).
    \item In the static case, the implied chained break-even GDP growth rate is \(1.80\%\) annualized versus \(1.86\%\) annualized in the direct \(Y \rightarrow U\) regression.
    \item In the time-varying case, the mean chained break-even rate is \(1.79\%\) annualized versus \(1.86\%\) annualized for the direct estimate.
    \item The main difference is volatility:
    the direct annualized series has a standard deviation of about \(0.20\) percentage points,
    while the chained series has a standard deviation of about \(0.43\) percentage points.
    \item Interpretation:
    the chained route broadly validates the level of the direct estimate, but it amplifies time variation and is therefore best read as a consistency exercise rather than the headline measure.
\end{itemize}

\begin{figure}[H]
    \centering
    \includegraphics[width=0.83\textwidth]{../outputs/memo_direct_vs_chained.png}
    \caption{Direct versus chained break-even GDP growth.}
\end{figure}

\section{A Diagnostic Tool: Labour-Implied GDP Growth}
\begin{itemize}
    \item As an additional diagnostic, we invert the labour-market equations and ask what GDP growth path would be most consistent with the observed unemployment path and the observed employment path.
    \item Concretely, we take the estimated time-varying labour-market equations as given, treat \(U\) or \(E\) as observed, treat \(Y\) as latent, and impose a separate AR(2) law of motion for GDP growth.
    \item The GDP AR(2) is estimated with the same two pandemic pulse dummies used elsewhere, but the inverse filter itself does not force those dummies into the latent GDP path.
    \item This is not an alternative measure of official GDP. It is a model-based consistency tool: if the labour-market equations are informative, the implied GDP lines should co-move with actual GDP.
    \item The results are encouraging but not exact. The unemployment-implied GDP series has a correlation of about \(0.45\) with actual GDP growth, while the employment-implied GDP series has a correlation of about \(0.49\).
    \item In the latest quarter, 2025 Q4, actual GDP growth is about \(0.6\%\) annualized, compared with \(1.7\%\) annualized for the unemployment-implied series and \(1.0\%\) annualized for the employment-implied series.
    \item Practical interpretation:
    this tool is useful as a triangulation device. When the labour-implied GDP lines move well away from measured GDP, that flags either unusual labour-market adjustment or limits of the reduced-form mapping.
\end{itemize}

\begin{figure}[H]
    \centering
    \includegraphics[width=0.83\textwidth]{../outputs/memo_implied_gdp_tool.png}
    \caption{Actual GDP growth and the GDP growth rates implied by the unemployment and employment equations.}
\end{figure}

\section{Conclusion}
\begin{itemize}
    \item The cleanest headline result remains the direct \(Y \rightarrow U\) specification.
    \item On that benchmark, Swiss break-even GDP growth is close to \(1.9\%\) annualized on average and just above \(2.0\%\) annualized in the latest quarter.
    \item The preferred calibration implies meaningful but not excessive time variation.
    \item The employment-based specifications are useful corroborating evidence:
    they broadly support the level of the direct benchmark when chained together, but they are less stable as standalone break-even measures.
    \item Recommended communication line:
    ``A reduced-form time-varying benchmark suggests that Swiss GDP growth of around \(2\%\) annualized is currently needed to keep unemployment broadly stable.'' 
\end{itemize}

\appendix

\section{Data Construction Summary}
\begin{itemize}
    \item Employment:
    a blended quarterly employment series is constructed by averaging q/q growth from the seasonally adjusted household and firm surveys, cumulating the average growth rate into an index, and scaling levels to match the household survey average in 2019.
    \item Household survey:
    the raw household-employment series is seasonally adjusted back to the 1970s using a quarterly STL-based procedure and then chained to the official seasonally adjusted BFS series from 2010 onward.
    \item Unemployment:
    the raw monthly seasonally adjusted series is aggregated to quarterly averages.
    \item GDP:
    the memo uses seasonally, calendar-, and sport-adjusted real GDP from SECO and its q/q growth rate.
    \item Final analysis variables:
    \(E\) is q/q growth of the constructed employment series,
    \(U\) is the quarterly change in the unemployment rate,
    and \(Y\) is q/q GDP growth.
\end{itemize}

\begin{figure}[H]
    \centering
    \includegraphics[width=0.83\textwidth]{../outputs/memo_data_construction.png}
    \caption{Data construction: the project blends household and firm employment information while preserving the household-level anchor.}
\end{figure}

\section{Data Appendix: Key Choices from Notebook 02}
\begin{itemize}
    \item The project uses both the BFS household and firm employment concepts rather than choosing only one source.
    \item The long household series is back-extended via internal seasonal adjustment, benchmarked to the official overlap.
    \item The final blended employment measure is intended to preserve household-level scaling while smoothing quarter-to-quarter survey-specific noise.
    \item Quarterly aggregation choices matter:
    unemployment is quarter-average,
    while GDP and employment are used in quarterly growth-rate form.
\end{itemize}

\begin{figure}[H]
    \centering
    \includegraphics[width=0.83\textwidth]{../outputs/memo_analysis_dataset.png}
    \caption{Final quarterly analysis dataset used in estimation: employment growth (\(E\)), change in unemployment (\(U\)), and GDP growth (\(Y\)).}
\end{figure}

\section{Empirical Caveats}
\begin{itemize}
    \item These are reduced-form relationships, not structural labour-market models.
    \item Results depend on the choice of calibration for the state variance; the memo uses \(\sigma_{\eta}=0.05\,\sigma_{OLS}\) as a deliberately smooth benchmark.
    \item The COVID pulse dummies improve robustness, but unusual episodes can still influence the estimated latent intercept.
    \item The \(Y \rightarrow E\) specification gives a different break-even interpretation from the direct \(Y \rightarrow U\) benchmark and should not be over-interpreted on its own.
    \item The chained \(Y \rightarrow E \rightarrow U\) comparison is an internal consistency check, not an independent identification strategy.
\end{itemize}

\end{document}
